\begin{abstract}
Token-level differentially private (DP) prompt optimization methods such as DP-OPT can become unstable under tight privacy budgets: on GSM8K, DP-OPT obtains $49.5{\pm}28.5\%$ across 30 runs, and a logged search trajectory reveals prompt-template drift and noise-sensitive irreversible choices. We diagnose these as structural consequences of greedy token-by-token construction over privately aggregated counts. We then propose \textbf{DP-ES} (Differentially Private Evolution Strategies), a structurally cleaner alternative that maintains a population of full prompts, mutates them via LLM calls that never access the private dataset, and spends privacy only on sampled-Gaussian evaluation; deterministic or Gumbel-smoothed selection is post-processing. Under a conservative $(\varepsilon{\leq}1.0, \delta{=}10^{-5})$ guarantee, DP-ES achieves \textbf{88.1\% on GSM8K} (+38.6 pp over DP-OPT, $\sim$9$\times$ lower std), \textbf{99.7\% on MedQA}, \textbf{73.5\% on BANKING77}, and \textbf{86.8\% on Alpaca}. It is also 2.5$\times$ faster in wall-clock time and uses 3.3$\times$ fewer logged private-data call groups than DP-OPT. Selection and population ablations, implementation-level noise checks, and a 200-profile exact-match memorization stress test complement the formal guarantee. \textbf{Scope.} Our experiments establish optimization robustness under DP noise, especially where prompt structure is critical; end-to-end validation on genuinely sensitive, non-saturated deployment data remains future work. Code: \href{https://github.com/StephCpa/dp-es}{github.com/StephCpa/dp-es}.
\end{abstract}
